%! AP Statistics — Unit 2, Lesson 2.7 — Homework — Answer Key
\documentclass[10pt]{article}
\usepackage{apstats-article}
\usepackage{apstats-key}

\begin{document}

\pageheader{Unit 2, Lesson 2.7}{Homework --- Answer Key}
\namedateperiod

\vspace{0.18in}

% ── DIRECTIONS ───────────────────────────────────────────────────────────────
\noindent\textbf{Directions:} Show all work. Write all interpretations in complete sentences
with context. Use variable names, not just $x$ and $y$.

\vspace{0.14in}

% ── PROBLEM 1 ─────────────────────────────────────────────────────────────────
\noindent\textbf{1.}\enspace A sports scientist tracked weekly training distance (miles, $x$)
and 5K race time (minutes, $y$) for 6 runners. Technology produced the regression equation
$\widehat{\text{time}} = 36 - 0.4(\text{miles})$. The table below gives the data.

\vspace{8pt}
\renewcommand{\arraystretch}{1.9}
\begin{tabularx}{\linewidth}{@{} >{\centering\arraybackslash}p{2.8cm}
                                    >{\centering\arraybackslash}p{2.2cm}
                                    >{\centering\arraybackslash}p{2.8cm}
                                    X @{}}
\rowcolor{sky}
\textbf{Training (mi, $x$)} & \textbf{Time (min, $y$)} &
\textbf{Predicted ($\hat{y}$)} & \textbf{Residual ($e$)} \\
\midrule
10 & 33 & \ans{32} & \ans{$+1$} \\
15 & 31 & \ans{30} & \ans{$+1$} \\
20 & 29 & \ans{28} & \ans{$+1$} \\
25 & 24 & \ans{26} & \ans{$-2$} \\
30 & 25 & \ans{24} & \ans{$+1$} \\
35 & 22 & \ans{22} & \ans{$0$} \\
\midrule
\multicolumn{3}{r}{\textbf{Sum of residuals}} & \ans{$+2$} \\
\end{tabularx}

\begin{teachernote}
The residuals sum to $+2$, not 0, because the given equation is not the true LSRL for this
dataset. Students should notice the discrepancy and raise it. Use this as a teachable moment:
only the actual LSRL guarantees $\sum e = 0$.
\end{teachernote}

\vspace{10pt}

\begin{enumerate}[label=\textbf{(\alph*)}, leftmargin=2em, itemsep=0.18in]

  \item Compute $\hat{y}$ and the residual for each runner. (See table above.)\\
        \ansline{Note: $\sum e = +2$ here, not 0, because the given equation is not the exact
        LSRL for these six data points.}

  \item Interpret the residual for the runner who trains \textbf{25 miles per week}.\\[8pt]
        \ansline{The residual is $e = 24 - 26 = -2$ minutes. The actual 5K time for the runner
        who trains 25 miles per week is 2 minutes faster than the value predicted by the
        regression line (the model overpredicted this runner's time).}

  \item Which runner has the largest absolute residual? What does this tell you?\\[8pt]
        \ansline{The 25-miles-per-week runner has the largest absolute residual ($|{-2}| = 2$
        min). This runner's 5K time is furthest from the regression line's prediction --- the
        model fits this data point less well than the others.}

\end{enumerate}

\vspace{0.18in}

% ── PROBLEM 2 ─────────────────────────────────────────────────────────────────
\noindent\textbf{2.}\enspace Two researchers fit linear regression models to different datasets.
Their residual plots are described below.

\vspace{8pt}
\begin{tabularx}{\linewidth}{@{} >{\bfseries}p{2.2cm} X @{}}
\rowcolor{sky} \textbf{Model} & \textbf{Description of residual plot} \\
\midrule
Model A & Residuals appear randomly scattered above and below zero with no obvious pattern. \\[4pt]
Model B & Residuals are negative in the middle of the $x$-range and positive at the extremes,
          forming a clear U-shape. \\
\end{tabularx}

\vspace{10pt}

\begin{enumerate}[label=\textbf{(\alph*)}, leftmargin=2em, itemsep=0.18in]

  \item Write a complete AP-style conclusion about \textbf{Model A}.\\[8pt]
        \ansline{The residual plot for Model A shows no apparent pattern (random scatter around
        zero). This suggests that a linear model is an appropriate model for the data.}

  \item Write a complete AP-style conclusion about \textbf{Model B}.\\[8pt]
        \ansline{The residual plot for Model B shows a clear curved (U-shaped) pattern.
        This suggests that a linear model is NOT an appropriate model for the data; a
        nonlinear model may provide a better fit.}

\end{enumerate}

\vspace{0.18in}

% ── PROBLEM 3 ─────────────────────────────────────────────────────────────────
\noindent\textbf{3.}\enspace A marine biologist uses water temperature (°C, $x$) to predict
coral growth rate (mm/month, $y$). Technology output from the regression of growth rate on
temperature is shown below.

\vspace{8pt}
\renewcommand{\arraystretch}{1.4}
\begin{tabularx}{\linewidth}{@{} l X X X X @{}}
\rowcolor{sky}
\textbf{Term} & \textbf{Coef} & \textbf{SE Coef} & \textbf{T-Value} & \textbf{P-Value} \\
\midrule
Constant       & $-0.85$ & 0.22 & $-3.86$ & 0.004 \\
Temperature    & $0.18$  & 0.01 &  18.0  & 0.000 \\
\end{tabularx}

\vspace{10pt}

\begin{enumerate}[label=\textbf{(\alph*)}, leftmargin=2em, itemsep=0.18in]

  \item Write the regression equation using variable names.\\[8pt]
        \ansline{$\widehat{\text{growth rate}} = -0.85 + 0.18(\text{temperature})$}

  \item A water sample at 24°C produced an actual growth rate of 3.9 mm/month. Calculate
        and interpret the residual in context.\\[8pt]
        \ansline{$\hat{y} = -0.85 + 0.18(24) = -0.85 + 4.32 = 3.47$ mm/month}\\[3pt]
        \ansline{$e = 3.9 - 3.47 = +0.43$ mm/month. The actual coral growth rate at 24°C is
        0.43 mm/month higher than the value predicted by the regression line.}

  \item The residual plot shows a curved pattern. Write a complete AP-style statement about
        whether a linear model is appropriate for these data.\\[8pt]
        \ansline{The residual plot shows a curved pattern. This suggests that a linear model
        may not be an appropriate model for predicting coral growth rate from water temperature;
        a nonlinear model may provide a better fit.}

\end{enumerate}

\vspace{0.18in}

% ── EXTENSION ─────────────────────────────────────────────────────────────────
\begin{extensionbox}
\small\textbf{Extension --- Optional}

For a dataset you find online (e.g., in a textbook, scientific article, or data repository),
fit a linear regression and create a residual plot using technology (calculator RESID list
or Desmos). Write: (1) the regression equation with variable names, (2) the largest
residual and its interpretation in context, and (3) your conclusion about whether a linear
model is appropriate based on the residual plot. Attach a screenshot or sketch.
\end{extensionbox}

\vspace{0.16in}

% ── PREVIEW ───────────────────────────────────────────────────────────────────
\begin{tcolorbox}[colback=greenbg, colframe=greenacc, boxrule=0.8pt, arc=2mm,
    left=4mm, right=4mm, top=3mm, bottom=3mm]
\small\textbf{Preview --- Lesson 2.8}\\[4pt]
The least-squares regression line minimizes $\sum e_i^2$ --- the sum of \textit{squared}
residuals. But why square them? And where do the formulas
$b = r \cdot \dfrac{s_y}{s_x}$ and $a = \bar{y} - b\bar{x}$
actually come from? In Lesson 2.8 we derive them and see how $r$, $s_x$, $s_y$, $\bar{x}$,
and $\bar{y}$ all work together to pin down the line that fits best.\\[6pt]
\textbf{Think ahead:} Why is minimizing $\sum e_i^2$ a better goal than minimizing $\sum e_i$?\\[6pt]
\ansline{Because $\sum e_i = 0$ for any least-squares line (positives and negatives cancel),
so minimizing the raw sum tells us nothing. Squaring the residuals makes all contributions
positive so the minimization actually penalizes large errors.}
\end{tcolorbox}

\end{document}
